\section{Реализация}
\label{sec:implementation}

В предыдущих главах были описаны теоретическая модель адаптивного
кодирования (разд.~\ref{sec:theory}) и архитектура решения, использующего эту
модель. В данной главе фиксируется конкретное воплощение этой архитектуры:
выбранные параметры, отображение архитектурных сущностей в структуры данных,
набор реализованных переходов перекодировщика, код проверочных блоков, модель
задержки чтения и перечень величин, выставляемых наружу для наблюдения.

\subsection{Параметры системы}

Параметры в табл.~\ref{tab:implementation-parameters} задают конкретный
вариант системы. Все значения масштабируются монотонно: переход к большему
числу узлов или большему лимиту диска расширяет пространство допустимых схем,
но не меняет ни параметрическое условие целевой отказоустойчивости
$r + g + \mathbf{1}_{l>0} = 3$, ни структуру переходов.

\begin{center}
\begin{longtable}{|p{0.23\textwidth}|p{0.25\textwidth}|p{0.42\textwidth}|}
\caption{\label{tab:implementation-parameters}
Основные параметры системы и экспериментального стенда.}\\
\hline
Параметр & Значение & Роль в системе \\
\hline
Зона доступности & одна зона доступности & зона размещения виртуальных машин стенда; отказ зоны не моделировался \\
\hline
Узлы хранения & instance group из \(30\) виртуальных машин & физическая основа стенда; каждая виртуальная машина соответствует одному домену отказа \\
\hline
Параметры VM & \(2\) vCPU, \(2\) ГБ RAM & вычислительные параметры виртуальных машин, на которых запускались контейнеры сервисов \\
\hline
Дополнительный диск VM & secondary disk \(13\) ГБ & диск виртуальной машины, используемый для размещения данных контейнеров стенда \\
\hline
Лимит ёмкости контейнера & \(1\) ГиБ на контейнер узла хранения & эффективная ёмкость узла в модели размещения; не является размером диска виртуальной машины \\
\hline
Размер блока & \(4\) МиБ & размер логического блока данных и единица учёта операций чтения \\
\hline
Базовая ширина полосы \(d_{\min}\) & \(6\) & минимальная допустимая ширина схемы и одновременно размер одного слота размещения \\
\hline
Число слотов размещения объекта & \(4\) & глубина дерева слотов; выбрано из условия вместимости: \(3\) (реплики) \(+\) \(3\) (локализованные узлы) \(+\) \(4 \cdot 6\) (данные в слотах) \(= 30\) \\
\hline
Локализованные узлы избыточности объекта на схему & \(3\) & узлы, выделяемые под проверочные блоки; число следует из условия \(r + g + \mathbf{1}_{l>0} = 3\): при \(l=1, g=2\) требуется \(3\) узла \\
\hline
Фиксированная широкая схема & \(\begin{array}{c}Hy(0,\\ LRCC(24,1,2))\end{array}\) & базовая сравнительная точка по плотности хранения \\
\hline
Период полураспада температуры & \(300\) секунд & постоянная времени экспоненциально взвешенного сглаживания (EWMA) температуры схемы \\
\hline
\end{longtable}
\end{center}

Инфраструктурный контур стенда не влияет на математическую модель кодирования,
но задаёт условия запуска прототипа: число доменов отказа, способ размещения
сервисов, источник контейнерных образов и путь сбора метрик. Это позволяет отделить
параметры экспериментальной среды от параметров самих схем
\(Hy(r,LRCC(d,l,g))\).

\begin{center}
\normalsize
\begin{tikzpicture}[
    infra/.style={
        draw, rounded corners, align=center,
        text width=3.2cm, minimum height=1.15cm,
        fill=orange!12, font=\normalsize
    },
    vm/.style={
        draw, rounded corners, align=center,
        text width=3.35cm, minimum height=1.25cm,
        fill=blue!7, font=\normalsize
    },
    svc/.style={
        draw, rounded corners, align=center,
        text width=3.45cm, minimum height=1.25cm,
        fill=green!8, font=\normalsize
    },
    aux/.style={
        draw, rounded corners, align=center,
        text width=3.25cm, minimum height=1.15cm,
        fill=purple!8, font=\normalsize
    },
    mon/.style={
        draw, rounded corners, align=center,
        text width=3.25cm, minimum height=1.15cm,
        fill=gray!10, font=\normalsize
    },
    edge/.style={-{Latex[length=2mm]}, thick},
    runedge/.style={-{Latex[length=2mm]}, dashed, thick},
    metricedge/.style={-{Latex[length=2mm]}, dotted, thick}
]

% 1. Infrastructure layer
\node[infra] (sa) at (-4.2,12.0) {Сервисный аккаунт\\\texttt{ig-sa}};
\node[infra] (subnet) at (4.2,12.0) {Подсеть\\ru-central1-a\\1 зона доступности};

% 2. VM layer
\node[vm] (ig) at (-6.3,8.4) {Instance Group\\30 VM};
\node[vm] (tpl) at (-2.1,8.4) {Шаблон VM\\2 vCPU, 2 ГБ RAM\\preemptible};
\node[vm] (nodes) at (2.1,8.4) {VM 1..30\\узлы хранения\\1 VM = 1 домен отказа};
\node[aux] (disk) at (6.3,8.4) {Secondary disk\\13 ГБ};

% 3. Runtime layer
\node[aux] (limit) at (-6.3,4.8) {Лимит ёмкости\\1 ГиБ на контейнер\\узла хранения};
\node[svc] (storage_services) at (-1.9,4.8) {Контейнеры сервисов\\на VM 1..30};
\node[infra] (reg) at (2.5,4.8) {Container Registry\\образы сервисов};

% 4. Metrics layer
\node[mon] (monium) at (0.3,1.4) {Monium\\сбор метрик};

% Infrastructure
\draw[edge] (sa) -- (ig);
\draw[edge] (ig) -- (tpl);
\draw[edge] (tpl) -- (nodes);
\draw[edge] (subnet) -- (nodes);
\draw[edge] (nodes) -- (disk);

% Runtime / startup
\draw[runedge] (nodes.south) -- (storage_services.north);
\draw[runedge] (reg.west) -- (storage_services.east);
\draw[runedge] (limit.east) -- (storage_services.west);

% Metrics
\draw[metricedge] (storage_services.south) -- (monium.north);

\end{tikzpicture}

\captionof{figure}{\label{fig:yc-terraform-stand}
Инфраструктурная схема экспериментального стенда в Yandex Cloud.
Сплошные стрелки на рис.~\ref{fig:yc-terraform-stand} показывают ресурсы,
создаваемые Terraform, пунктирные --- запуск контейнеров и взаимодействие
сервисов, точечные --- сбор метрик.
}
\end{center}

В этом варианте абстрактные вероятности отказов из разд.~\ref{sec:theory-read-cost}
превращаются в конкретные числа. Для \(H = 30\) узлов и \(K = 3\)
недоступных узлов гипергеометрические вероятности
\(p_0^*, p_1^*, p_2^*, p_{3+}\) вычисляются заранее, а ожидаемая стоимость
чтения \(E[io]\) используется планировщиком как априорная оценка стоимости
каждой допустимой схемы вместе с температурой этой схемы.

\subsection{Структуры данных}

Архитектурные сущности из разд.~\ref{sec:theory-hybrid} получают конкретное
представление в виде отдельных структур данных
(табл.~\ref{tab:implementation-classes}). Физические блоки проиндексированы
на узлах, а каждый блок адресуется парой (узел, локальный ключ). Это устраняет
необходимость в глобальном именном индексе блоков и делает физический слой
плоским.

\begin{center}
\begin{longtable}{|p{0.25\textwidth}|p{0.65\textwidth}|}
\caption{\label{tab:implementation-classes}
Основные структуры данных системы.}\\
\hline
Сущность & За что отвечает \\
\hline
\texttt{ChunkRef} & ссылка на блок: узел, локальный ключ \\
\hline
\texttt{ReplicaLayout} & упорядоченный список ссылок на физические блоки реплики \\
\hline
\texttt{LrccLayout} & LRCC-представление полосы: блоки данных, локальные и глобальные проверочные блоки \\
\hline
\texttt{Scheme} & состояние избыточности схемы $Hy(r, LRCC(d, l, g))$, температура и размещение \\
\hline
\texttt{ObjectExtent} & логический диапазон объекта, связанный с конкретной схемой и подмножеством используемых слотов \\
\hline
\texttt{ObjectRecord} & метаданные объекта: размер, последовательность диапазонов, набор слотов, дерево совместимости слотов, выбранные узлы для локализованного хранения блоков избыточности \\
\hline
\texttt{Slot} & группа узлов, по которым раскладываются последовательные блоки данных схемы; единица проверки совместимости при объединении/разделении \\
\hline
\texttt{SlotTree} & структура, задающая правило дерева слотов для допустимых объединений и разделений \\
\hline
\end{longtable}
\end{center}

Порядок блоков данных в списках \texttt{ReplicaLayout} и \texttt{LrccLayout}
совпадает с логическим порядком байтов. Благодаря этому для чтения логического
диапазона достаточно вычислить смещение внутри \texttt{ObjectExtent}, не
строя отдельных индексов по физическим блокам.

\subsection{Кодирование проверочных блоков}

Локальный проверочный блок строится как побитовое исключающее ИЛИ (XOR) всех
блоков данных одной локальной группы. Это даёт максимально дешёвый по
вычислительной сложности способ восстановить любой один потерянный блок данных
внутри своей группы.

Глобальные проверочные блоки строятся как линейные комбинации всех блоков
данных полосы по строкам матрицы Вандермонда над полем $GF(2^8)$, как в кодах
Рида--Соломона \cite{reed1960polynomial}. Для объединения и разделения полос
используется идея конвертируемых кодов
\cite{maturana2022convertible,kong2024lrcconvertible}: эти переходы выполняются
без перечитывания блоков данных --- только за счёт линейных операций над уже
имеющимися проверочными блоками. В частности, новые глобальные проверочные
блоки получаются линейными преобразованиями существующих проверочных блоков и
не требуют обращения ко всем блокам данных полосы.

Восстановление при многократных отказах выполняется общим решателем линейных
уравнений над $GF(2^8)$. Он собирает доступные блоки данных и доступные
локальные и глобальные проверочные блоки, строит систему уравнений с числом
уравнений не меньше числа неизвестных и решает её для пропавших блоков. В тех
состояниях \(Hy(r,LRCC(d,l,g))\), которые удовлетворяют условию
\(r + g + \mathbf{1}_{l>0}=3\), такой решатель позволяет обслуживать
пользовательское чтение при отказах до трёх доменов при условии разнесения
блоков по доменам отказа и достаточного числа независимых уравнений.

\subsection{Поддержанные переходы перекодировщика}

В архитектурной модели (разд.~\ref{sec:theory-hybrid}) пространство допустимых
переходов между схемами задаётся параметрически. В моделируемой системе из
этого пространства реализованы переходы, перечисленные в
табл.~\ref{tab:implementation-transitions}. Каждое направление имеет свой
источник данных и свой набор удаляемых после фиксации блоков.

Составные переходы обмена реплик на проверочные блоки (например, переход от
двух дополнительных реплик к набору из одного локального и двух глобальных
проверочных блоков, описанный в нижней строке таблицы) поддержаны явно: блоки
данных читаются ровно один раз, а затем из них вычисляются все требуемые
проверочные блоки. Это исключает промежуточное состояние и предотвращает
многократную оплату чтения данных при создании нескольких новых проверочных
блоков.

\begin{center}
\begin{longtable}{|p{0.30\textwidth}|p{0.24\textwidth}|p{0.36\textwidth}|}
\caption{\label{tab:implementation-transitions}
Поддержанные физические переходы перекодировщика.}\\
\hline
Тип перехода & Направление & Что физически происходит \\
\hline
Реплика в локальный проверочный блок & к более компактной схеме & создаётся один локальный проверочный блок (XOR блоков данных локальной группы), удаляется последняя реплика \\
\hline
Реплики в глобальные проверочные блоки & к более компактной схеме & создаются один или два глобальных проверочных блока, удаляются соответствующие реплики \\
\hline
Все дополнительные реплики во все проверочные блоки & к LRCC-представлению & создаются один локальный и два глобальных проверочных блока, удаляются все дополнительные реплики, остаётся только LRCC-представление полосы \\
\hline
Проверочные блоки в реплики & к более дешёвому чтению & создаются новые реплики, удаляются выбранные проверочные блоки \\
\hline
Разделение локальных групп ($l \to 2l$) & к более локальному восстановлению & пересчитываются новые локальные проверочные блоки для меньших групп \\
\hline
Объединение локальных групп ($l \to l/2$) & к более компактной схеме & новые локальные проверочные блоки строятся из старых их объединением \\
\hline
Объединение схем (объединение полос) & к более широкой схеме & две схемы одного объекта, совместимые по дереву слотов, получают общую схему $Hy(r, LRCC(2d,2l,g))$, после чего новые глобальные проверочные блоки получаются линейным преобразованием уже имеющихся проверочных блоков без перечитывания данных \\
\hline
Разделение схемы (разделение полосы) & к двум более узким схемам & общая схема разбивается на две схемы с исходными подмножествами занятых слотов; глобальные проверочные блоки дочерних схем также получаются линейным преобразованием уже имеющихся проверочных блоков без перечитывания данных \\
\hline
\end{longtable}
\end{center}

При создании новой реплики система предпочитает копировать существующую реплику,
если хотя бы одна осталась; только если в схеме реплик уже нет, исходные блоки
данных читаются из LRCC-представления. Эта асимметрия влияет на оценку
стоимости нагрева: схема, полностью ушедшая в кодовое представление, при
необходимости вернуть в неё реплики обойдётся дороже.

\subsection{Алгоритмы планировщиков}
\label{sec:implementation-planners}

В прототипе локальный планировщик за один проход сохраняет только один лучший
кандидат. Этот кандидат может соответствовать одиночному переходу или точной
паре объединения, найденной на основе совместимости слотов.

\begin{algorithm}[H]
\SetAlgoLined
\KwData{виртуальные\_метрики}
\KwResult{лучший\_переход}
лучший\_переход $\leftarrow$ нет\;
индекс\_кандидатов\_объединения $\leftarrow$ пусто\;

\For{каждая схема}{
    одиночный\_переход $\leftarrow$ лучший\_одиночный\_переход(схема)\;
    лучший\_переход $\leftarrow$ max(лучший\_переход, одиночный\_переход)\;

    \If{одиночный\_переход ведёт к объединению}{
        профиль $\leftarrow$ (объект, текущая\_сигнатура, целевая\_сигнатура)\;
        обновить\_максимумы(индекс\_кандидатов\_объединения, профиль, схема)\;
    }
}

\For{каждый (профиль, маска) из индекс\_кандидатов\_объединения}{
    пара $\leftarrow$ найти\_пару\_маски\_по\_дереву\_совместимости(индекс\_кандидатов\_объединения, профиль, маска)\;
    лучший\_переход $\leftarrow$ max(лучший\_переход, пара)\;
}

сохранить лучший\_переход\;

\caption{Упрощённый псевдокод локального планировщика}
\end{algorithm}

Глобальный планировщик допускает переходы без ожидания завершения
перекодировщика: после каждого допуска он сразу применяет ожидаемый эффект к
виртуальным метрикам. Учёт допущенных, но ещё не завершённых переходов
остаётся на стороне локального планировщика как состояние переходов в
исполнении.

\begin{algorithm}[H]
\SetAlgoLined
\KwData{лимит}
\KwResult{допущенные\_переходы}
допущенные\_переходы $\leftarrow$ пусто\;
виртуальные\_метрики $\leftarrow$ текущие\_метрики\_кластера()\;

\While{есть лимит фоновых переходов}{
    кандидаты $\leftarrow$ опросить\_локальные\_планировщики(виртуальные\_метрики)\;
    кандидат $\leftarrow$ лучший\_положительный(кандидаты)\;

    \If{кандидат отсутствует}{
        продолжить\;
    }

    выдать\_допуск(ответственный\_локальный\_планировщик, кандидат)\;
    добавить(допущенные\_переходы, кандидат)\;
    виртуальные\_метрики $\leftarrow$ применить\_эффект(виртуальные\_метрики, кандидат)\;
    продолжить без ожидания перекодировщика\;
}

вернуть допущенные\_переходы\;
\caption{Упрощённый псевдокод глобального планировщика}
\end{algorithm}

Между принятием решения о переходе и его фактическим завершением проходит
некоторое время. Кандидат, получивший разрешение глобального планировщика,
передаётся локальному планировщику для запуска перекодирования. Пока переход
находится в исполнении, соответствующий планировщик считается занятым и для
него не выбираются новые переходы. Это снижает вероятность слишком резких
разовых изменений и ограничивает нагрузку на узлы, на которых локализованы
проверочные блоки одного объекта.

\subsection{Модель чтения данных из разных схем}

Задержка одной операции чтения на узле описывается суммой накладных расходов
базовой задержки и времени передачи, пропорционального объёму реально
прочитанных данных:
\[
    t_{read}(b) = t_{base} + \frac{b}{B},
\]
где $t_{base}$ --- базовая задержка узла, $b$ --- объём прочитанных данных,
$B$ --- пропускная способность канала ввода-вывода.

Клиент читает данные одним из способов:

\begin{itemize}
  \item для чтения из реплики предполагается потоковая выдача из самой быстрой
  реплики в схеме. Диапазон, затрагивающий несколько последовательных
  логических блоков, читается как один последовательный поток с одного узла.
  Тем самым несколько блоков могут быть получены за одну физическую операцию
  чтения, а логические блоки становятся доступными клиенту по мере чтения;
  \item для чтения из LRCC предполагается сначала прочитать все необходимые
  блоки данных: запрос диапазона, затрагивающий $k$ блоков одной полосы,
  требует дождаться всех $k$ чтений, которые можно выполнять параллельно;
  соответственно задержка будет максимальной из соответствующих узлов в
  схеме. Далее данные собираются в исходном порядке: выполняется склейка
  прочитанных блоков и, при необходимости, чтение с восстановлением. В отличие
  от реплик, соседние блоки в LRCC чаще всего размещены на разных узлах,
  поэтому даже последовательный диапазон обычно требует нескольких чтений.
\end{itemize}

В режиме чтения с восстановлением к $k$ добавляются чтения вспомогательных и
проверочных блоков, требуемых для решения системы линейных уравнений либо
локального восстановления.

\subsection{Метрики мониторинга}

Система поставляет следующие метрики, по которым можно судить о состоянии
кластера. Их структура приведена в табл.~\ref{tab:monitoring-metrics}.
Имена метрик и типы \texttt{gauge}, \texttt{counter}, \texttt{histogram} и
\texttt{summary} оставлены в форме, принятой в системах мониторинга.

\newcommand{\metriccell}[2]{%
\begin{tabular}[t]{@{}l@{}}
\texttt{#1}\\
\texttt{#2}
\end{tabular}}

\providecommand{\metriccell}[2]{%
\begin{tabular}[t]{@{}l@{}}
\texttt{#1}\\
\texttt{#2}
\end{tabular}}

\begin{center}
\small
\begin{longtable}{|p{0.39\textwidth}|p{0.16\textwidth}|p{0.37\textwidth}|}
\caption{\label{tab:monitoring-metrics}
Метрики мониторинга прототипа.}\\
\hline
Метрика & Тип и метки & Значение \\
\hline
\endfirsthead

\hline
Метрика & Тип и метки & Значение \\
\hline
\endhead

\texttt{storage\_used\_bytes}
&
\texttt{gauge}

Метки: \texttt{kind}
&
Занятое место по типам блоков: данные, реплики, локальные и глобальные
проверочные блоки.

Компонент: хранилище / узлы данных. \\
\hline

\texttt{allocator\_node\_free\_bytes}
&
\texttt{gauge}

Метки: \texttt{node}
&
Свободное место на узле для новых блоков.

Компонент: распределитель. \\
\hline

\texttt{scheme\_distribution}
&
\texttt{gauge}

Метки: \texttt{r,d,l,g}
&
Число физических записей \texttt{Scheme} с данной сигнатурой.

Компонент: глобальный планировщик / метаданные. \\
\hline

\metriccell{logical\_chunks}{\_distribution}
&
\texttt{gauge}

Метки: \texttt{r,d,l,g}
&
Число логических блоков или байтов пользовательских данных.

Компонент: глобальный планировщик / метаданные. \\
\hline

\texttt{extent\_temperature}
&
\texttt{gauge}

Метки: \texttt{object, extent}
&
EWMA-интенсивность чтения диапазона или схемы.

Компонент: локальный планировщик. \\
\hline

\metriccell{read\_ops}{\_total}
&
\texttt{counter}

Метки: \texttt{mode}
&
Накопительное число физических чтений блоков с узлов, включая чтения данных и
проверочных блоков при восстановлении.

Компонент: модуль чтения. \\
\hline

\texttt{read\_latency\_ms}
&
\texttt{histogram} или \texttt{summary}

Метки: \texttt{mode}
&
Задержка чтения в миллисекундах: обычный режим / режим частичной недоступности.

Компонент: модуль чтения. \\
\hline

\metriccell{recoder\_transitions}{\_total}
&
\texttt{counter}

Метки: \texttt{action, result}
&
Число завершённых, отклонённых или неуспешных переходов.

Компонент: перекодировщик. \\
\hline

\texttt{recoder\_in\_flight}
&
\texttt{gauge}

Метки: \texttt{action}
&
Число допущенных, но ещё не завершённых переходов.

Компонент: локальный планировщик. \\
\hline

\metriccell{read\_logical\_chunks}{\_total}
&
\texttt{counter}

Метки: \texttt{mode}
&
Число логических блоков, запрошенных пользовательскими операциями чтения.
Используется как знаменатель при вычислении числа физических чтений на один
логический блок.

Компонент: модуль чтения. \\
\hline

\metriccell{recoder\_work\_bytes}{\_total}
&
\texttt{counter}

Метки: \texttt{action, kind}
&
Объём данных, прочитанных или записанных фоновым перекодировщиком.
Метка \texttt{kind} принимает значения чтения и записи, а \texttt{action}
задаёт тип перехода.

Компонент: перекодировщик. \\
\hline
\end{longtable}
\normalsize
\end{center}

Метрики распределения схем не являются гистограммами в смысле бакетов. Это
набор серий типа \texttt{gauge}, где точная сигнатура \(Hy(r, LRCC(d,l,g))\)
задаётся метками. Такое представление даёт возможность отслеживать не только
итоговый баланс «реплики против LRCC», но и качественные сдвиги внутри
LRCC-семейства, например миграцию от узких полос к широким.

\pagebreak
